\section{3. Background}

\subsection{3.1 Multi-Head Attention}

For input $X \in \mathbb{R}^{B \times L \times d}$:

\[
Q = X W_Q, \quad K = X W_K, \quad V = X W_V \in \mathbb{R}^{B \times L \times d}
\tag{3.1}
\]

Standard attention:

\[
\text{Attn}(X) = \text{softmax}\!\left(\frac{Q K^\top}{\sqrt{d_h}}\right) V
\tag{3.2}
\]

\subsection{3.2 K-Means Clustering}

Given $L$ points $x_1, \dots, x_L \in \mathbb{R}^d$, K-means minimizes:

\[
\mathcal{L}(C) = \sum_{i=1}^{L} \min_{k \in [K]} \|x_i - c_k\|_2^2
\tag{3.3}
\]

Lloyd's algorithm alternates:

1. \textbf{Assignment:} $a_i = \arg\min_k \|x_i - c_k\|_2^2$
2. \textbf{Update:} $c_k = \frac{\sum_{i: a_i = k} x_i}{|a_i = k|}$

\textbf{Complexity.} $\mathcal{O}(L K d)$ per iteration.

\textbf{Vectorization.} Both steps parallelize over $B, H$ in attention context.

\subsection{3.3 Differentiable K-Means}

Standard K-means uses \texttt{argmin}, which is non-differentiable. For 
end-to-end training:

\begin{itemize}
\item \textbf{Forward:} Use \texttt{argmin} to get hard assignments (for the update step).
\end{itemize}
\begin{itemize}
\item \textbf{Backward:} Use softmax for gradient flow through assignments.
\end{itemize}

ClusterFormer is fully differentiable:
\begin{itemize}
\item Centroid updates use one-hot (non-differentiable) — this is fine since
\end{itemize}
  centroids are deterministic functions of $Q$.
\begin{itemize}
\item Assignment uses \texttt{softmax} with learnable temperature, giving gradients.
\end{itemize}

\subsection{3.4 Linear Attention}

A generic linear attention reformulation:

\[
\text{Attn}(Q, K, V) = \frac{\phi(Q) \sum_j \phi(K_j) V_j^\top}{\phi(Q) \sum_j \phi(K_j)}
\tag{3.4}
\]

for some feature map $\phi$. If $\phi$ has dimension $m$, complexity is 
$\mathcal{O}(L m d)$.

ClusterFormer does not use this formulation. Instead, it clusters queries 
and computes attention from $K$ centroids to keys.

\subsection{3.5 LoRA Adapters}

LoRA (Hu et al., 2021) parameterizes:

\[
W_\text{eff} = W_\text{frozen} + \frac{\alpha}{r_L} B A, \quad B = 0 \text{ at init}
\tag{3.5}
\]

\subsection{3.6 Notation Summary}
